% GENERATED FILE. Edit canonical lessons, then run npm run generate:books.
% canonical-source-hash: fnv1a64:f2ef98f4548cb9a1
% canonical-lessons: TE-C358-jivincu, TE-C358-sthirapadu, TE-C358-taladacuko, TE-C358-kalaganu, TE-C358-adrstavasattu

\chapter{Making a Life}
\label{ch:te-making-a-life}
\hlchaptermodality{358}

\begin{chapteropening}
\textbf{By the end of this chapter:} \emph{I can say to live, to settle down, to take shelter, to dream of and luckily.}

Complete the last lesson of chapter 358: I can say to live, to settle down, to take shelter, to dream of and luckily.
\end{chapteropening}

\section[\texorpdfstring{\te{జీవించు} (jīviñcu)}{జీవించు (jīviñcu)}]{\te{జీవించు} (jīviñcu) — to live; to lead a life}

\label{lesson:TE-C358-jivincu}



\begin{quote}
Before the new one: say the Telugu for an orphan, then the Telugu for to survive.
\end{quote}

\subsection*{You'll want to know: \te{జీవించు}}
\textbf{\te{జీవించు}} — \emph{jīviñcu} — \textquotedblleft{}to live; to lead a life\textquotedblright{}.

\begin{quote}
\te{అతను} \te{వంద} \te{సంవత్సరాలు} \te{జీవించాడు}.
\end{quote}

\begin{quote}
\emph{atanu vanda saṁvatsarālu jīviñcāḍu.}
\end{quote}

\begin{quote}
He lived for a hundred years.
\end{quote}

Past (I …): \te{జీవించాను}, \emph{jīviñcānu}.

A formal partner of \te{బతుకు}.

\begin{cousinweb}
\textbf{\te{జీవించు}} (\emph{jīviñcu}) is built on Sanskrit \emph{jīvati}, \textquotedblleft{}is alive\textquotedblright{} (CDIAL 5241), with the Telugu verb ending \emph{-iñcu}.
\end{cousinweb}

\begin{grammarlens}[title={What you've built}]
One more word for telling what happened.
\end{grammarlens}

\subsection*{Your turn}
\begin{itemize}
\raggedright
  \item \emph{Say it:} \emph{jīviñcu}
  \item \emph{Say it:} \emph{jīviñcu}, once more
  \item \emph{Say it:} \emph{atanu vanda saṁvatsarālu jīviñcāḍu.} — He lived for a hundred years.
  \item \emph{Recall:} say the Telugu for an orphan, then the Telugu for to survive, then say \emph{jīviñcu} again
\end{itemize}

\begin{tcolorbox}[breakable,colback=teal!4,colframe=teal!35!black,title={Before you move on}]
What does \te{జీవించు} mean? (\textquotedblleft{}To live; to lead a life\textquotedblright{}.) Say it once more.
\end{tcolorbox}
\section[\texorpdfstring{\te{స్థిరపడు} (sthirapaḍu)}{స్థిరపడు (sthirapaḍu)}]{\te{స్థిరపడు} (sthirapaḍu) — to settle down}

\label{lesson:TE-C358-sthirapadu}



\begin{quote}
Before the new one: say the Telugu for to survive, then the Telugu for to live.
\end{quote}

\subsection*{You'll want to know: \te{స్థిరపడు}}
\textbf{\te{స్థిరపడు}} — \emph{sthirapaḍu} — \textquotedblleft{}to settle down\textquotedblright{}.

\begin{quote}
\te{చదువు} \te{తరువాత} \te{ఆమె} \te{పట్నంలో} \te{స్థిరపడింది}.
\end{quote}

\begin{quote}
\emph{caduvu taruvāta āme paṭnaṁlō sthirapaḍindi.}
\end{quote}

\begin{quote}
After her studies she settled in the city.
\end{quote}

Past (I …): \te{స్థిరపడ్డాను}, \emph{sthirapaḍḍānu}.

\te{స్థిరం}, firm or steady, with \te{పడు}.

\begin{cousinweb}
\textbf{\te{స్థిరపడు}} (\emph{sthirapaḍu}) is built on Sanskrit \emph{sthira}, \textquotedblleft{}firm, durable\textquotedblright{} (CDIAL 13771), with the Telugu helping verb \emph{paḍu}.
\end{cousinweb}

\begin{grammarlens}[title={What you've built}]
One more word for telling what happened.
\end{grammarlens}

\subsection*{Your turn}
\begin{itemize}
\raggedright
  \item \emph{Say it:} \emph{sthirapaḍu}
  \item \emph{Say it:} \emph{sthirapaḍu}, once more
  \item \emph{Say it:} \emph{caduvu taruvāta āme paṭnaṁlō sthirapaḍindi.} — After her studies she settled in the city.
  \item \emph{Recall:} say the Telugu for to survive, then the Telugu for to live, then say \emph{sthirapaḍu} again
\end{itemize}

\begin{tcolorbox}[breakable,colback=teal!4,colframe=teal!35!black,title={Before you move on}]
What does \te{స్థిరపడు} mean? (\textquotedblleft{}To settle down\textquotedblright{}.) Say it once more.
\end{tcolorbox}
\section[\texorpdfstring{\te{తలదాచుకో} (taladācukō)}{తలదాచుకో (taladācukō)}]{\te{తలదాచుకో} (taladācukō) — to take shelter}

\label{lesson:TE-C358-taladacuko}



\begin{quote}
Before the new one: say the Telugu for to live, then the Telugu for to settle down.
\end{quote}

\subsection*{You'll want to know: \te{తలదాచుకో}}
\textbf{\te{తలదాచుకో}} — \emph{taladācukō} — \textquotedblleft{}to take shelter\textquotedblright{}.

\begin{quote}
\te{వర్షం} \te{వచ్చి} \te{మేము} \te{ఒక} \te{గుడిలో} \te{తలదాచుకున్నాము}.
\end{quote}

\begin{quote}
\emph{varṣaṁ vacci mēmu oka guḍilō taladācukunnāmu.}
\end{quote}

\begin{quote}
It rained, and we took shelter in a temple.
\end{quote}

Past (I …): \te{తలదాచుకున్నాను}, \emph{taladācukunnānu}.

Literally to hide one's head: \te{తల}, head, and \te{దాచుకో}, to hide for oneself.

\begin{cousinweb}
\textbf{\te{తలదాచుకో}} (\emph{taladācukō}) is a native compound: \emph{tala}, \textquotedblleft{}the head\textquotedblright{} (DEDR 3103), and \emph{dācu}, \textquotedblleft{}to hide\textquotedblright{} (DEDR 63), with the self-ending \emph{-kō}.
\end{cousinweb}

\begin{grammarlens}[title={What you've built}]
One more word for telling what happened.
\end{grammarlens}

\subsection*{Your turn}
\begin{itemize}
\raggedright
  \item \emph{Say it:} \emph{taladācukō}
  \item \emph{Say it:} \emph{taladācukō}, once more
  \item \emph{Say it:} \emph{varṣaṁ vacci mēmu oka guḍilō taladācukunnāmu.} — It rained, and we took shelter in a temple.
  \item \emph{Recall:} say the Telugu for to live, then the Telugu for to settle down, then say \emph{taladācukō} again
\end{itemize}

\begin{tcolorbox}[breakable,colback=teal!4,colframe=teal!35!black,title={Before you move on}]
What does \te{తలదాచుకో} mean? (\textquotedblleft{}To take shelter\textquotedblright{}.) Say it once more.
\end{tcolorbox}
\section[\texorpdfstring{\te{కలగను} (kalaganu)}{కలగను (kalaganu)}]{\te{కలగను} (kalaganu) — to dream of}

\label{lesson:TE-C358-kalaganu}



\begin{quote}
Before the new one: say the Telugu for to settle down, then the Telugu for to take shelter.
\end{quote}

\subsection*{You'll want to know: \te{కలగను}}
\textbf{\te{కలగను}} — \emph{kalaganu} — \textquotedblleft{}to dream of\textquotedblright{}.

\begin{quote}
\te{చిన్నప్పుడు} \te{నేను} \te{డాక్టరు} \te{కావాలని} \te{కలగన్నాను}.
\end{quote}

\begin{quote}
\emph{cinnappuḍu nēnu ḍākṭaru kāvālani kalagannānu.}
\end{quote}

\begin{quote}
As a child I dreamed of becoming a doctor.
\end{quote}

Past (I …): \te{కలగన్నాను}, \emph{kalagannānu}.

\te{కల}, a dream, with \te{కను}, to see: to see a dream.

\begin{cousinweb}
\textbf{\te{కలగను}} (\emph{kalaganu}) is native Dravidian: DEDR 1407 lists Telugu \emph{kala}, \textquotedblleft{}a dream\textquotedblright{}, and \emph{kalaganu}, \textquotedblleft{}to dream\textquotedblright{}.
\end{cousinweb}

\begin{grammarlens}[title={What you've built}]
One more word for telling what happened.
\end{grammarlens}

\subsection*{Your turn}
\begin{itemize}
\raggedright
  \item \emph{Say it:} \emph{kalaganu}
  \item \emph{Say it:} \emph{kalaganu}, once more
  \item \emph{Say it:} \emph{cinnappuḍu nēnu ḍākṭaru kāvālani kalagannānu.} — As a child I dreamed of becoming a doctor.
  \item \emph{Recall:} say the Telugu for to settle down, then the Telugu for to take shelter, then say \emph{kalaganu} again
\end{itemize}

\begin{tcolorbox}[breakable,colback=teal!4,colframe=teal!35!black,title={Before you move on}]
What does \te{కలగను} mean? (\textquotedblleft{}To dream of\textquotedblright{}.) Say it once more.
\end{tcolorbox}
\section[\texorpdfstring{\te{అదృష్టవశాత్తు} (adṛṣṭavaśāttu)}{అదృష్టవశాత్తు (adṛṣṭavaśāttu)}]{\te{అదృష్టవశాత్తు} (adṛṣṭavaśāttu) — luckily; fortunately}

\label{lesson:TE-C358-adrstavasattu}



\begin{quote}
Before the new one: say the Telugu for to take shelter, then the Telugu for to dream of.
\end{quote}

\subsection*{You'll want to know: \te{అదృష్టవశాత్తు}}
\textbf{\te{అదృష్టవశాత్తు}} — \emph{adṛṣṭavaśāttu} — \textquotedblleft{}luckily; fortunately\textquotedblright{}.

\begin{quote}
\te{అదృష్టవశాత్తు} \te{ఆ} \te{ప్రమాదంలో} \te{ఎవరికీ} \te{ఏమీ} \te{కాలేదు}.
\end{quote}

\begin{quote}
\emph{adṛṣṭavaśāttu ā pramādaṁlō evarikī ēmī kālēdu.}
\end{quote}

\begin{quote}
Luckily nobody was hurt in that accident.
\end{quote}

\te{అదృష్టం}, luck, with \te{వశాత్తు}, by the power of.

\begin{cousinweb}
\textbf{\te{అదృష్టవశాత్తు}} (\emph{adṛṣṭavaśāttu}) is built from Sanskrit \emph{adṛṣṭa}, \textquotedblleft{}unseen, unforeseen\textquotedblright{} (Monier-Williams), and \emph{vaśāt}, \textquotedblleft{}by force of\textquotedblright{}, from \emph{vaśa}, \textquotedblleft{}will, power\textquotedblright{} (CDIAL 11430).
\end{cousinweb}

\begin{grammarlens}[title={What you've built}]
One more word for telling what happened.
\end{grammarlens}

\subsection*{Your turn}
\begin{itemize}
\raggedright
  \item \emph{Say it:} \emph{adṛṣṭavaśāttu}
  \item \emph{Say it:} \emph{adṛṣṭavaśāttu}, once more
  \item \emph{Say it:} \emph{adṛṣṭavaśāttu ā pramādaṁlō evarikī ēmī kālēdu.} — Luckily nobody was hurt in that accident.
  \item \emph{Recall:} say the Telugu for to take shelter, then the Telugu for to dream of, then say \emph{adṛṣṭavaśāttu} again
\end{itemize}

\begin{tcolorbox}[breakable,colback=teal!4,colframe=teal!35!black,title={Before you move on}]
What does \te{అదృష్టవశాత్తు} mean? (\textquotedblleft{}Luckily; fortunately\textquotedblright{}.) Say it once more.
\end{tcolorbox}
